% SSSUP-FONTE-STORICA: fonti-locali/originali/2014-2020.zip
% Anno 2016, esercizio 3. Testo storico; figura ricostruita da PDF.
% Verificare editorialmente la trasposizione prima di considerarla definitiva.
\renewcommand{\Anno}{2016}
\renewcommand{\NumeroProblema}{3}

\renewcommand{\Testo}{
La Scuola Sant'Anna vuole allestire una lotteria tra i propri studenti. Chi vuole giocare compra un biglietto e la sua vincita viene determinata nella maniera seguente: vengono posizionate due urne uguali, ognuna delle quali contiene 10 palline, numerate da 0 a 9, il giocatore estrae a caso una pallina da ognuna delle due urne e riceverà una vincita in Euro pari al prodotto dei numeri estratti.

\begin{enumerate}
\item[(a)] Determinare il prezzo minimo del biglietto (in Euro) per cui in media la Scuola non andrà in perdita allestendo la lotteria.
\item[(b)] Approfittando della distrazione dei commissari, uno studente sposta per scherzo una pallina, presa a caso, da un'urna all'altra; stabilire il prezzo minimo del biglietto, come sopra, in questa nuova situazione.
\end{enumerate}
}

\renewcommand{\Risposta}{}
\renewcommand{\Aiutino}{}
\renewcommand{\SoluzioneUfficiale}{}
\renewcommand{\SoluzioneRedazionale}{}
\renewcommand{\SoluzioniAlternative}{}
\renewcommand{\Osservazioni}{}
\renewcommand{\Commenti}{}
